% Volume I, appendice «Elementi di statica grafica»: figura fig:cremoniano_sgr.
% Richiamata con \FiguraSGR{figura-sgr-cremoniano}; didascalia e label nell'appendice.
% Generata con gen_sgr.py (coordinate numeriche; stili di Preambolo/tikz-pol e tikz-sgr).
%
% Capriata palladiana dell'esempio ese:capriata_palladiana (capitolo Sistemi reticolari) e suo diagramma
% cremoniano. Nodi (in l): A(0,0), B(2,0), C(4,0), D(1,1), E(3,1), F(2,2); forze F in D, F, E; reazioni 3F/2.
% Campi esterni 0, 1, 2, 3, m numerati in senso orario; campi interni I (ADB), II (DFB), III (FEB), IV (ECB).
% Vertici (F = 1): a = (-1.500,-1.500), b = (-1.000,-2.000), c = (-1.000,-1.000), d = (-1.500,-1.500) (a e d coincidono, per simmetria).
% Sforzi letti sul diagramma, uguali a quelli dell'esempio (verificato): AB = BC = 3/2, BF = 1 (tiranti);
% AD = CE = 3 sqrt2/2, DF = FE = sqrt2, DB = EB = sqrt2/2 (puntoni).
% Capriata in scala 1.3 cm per l; cremoniano in scala 1.7 cm per F.
% (a) capriata liberata dai vincoli, con carichi, reazioni e campi; (b) diagramma cremoniano: tratto spesso
%     per le aste compresse, sottile per le tese.
\begin{tikzpicture}[font=\footnotesize]
\begin{scope}[x=1.3cm,y=1.3cm]
\draw[sgrasta] (0.000,0.000) -- (1.000,1.000);
\draw[sgrasta] (0.000,0.000) -- (2.000,0.000);
\draw[sgrasta] (1.000,1.000) -- (2.000,2.000);
\draw[sgrasta] (1.000,1.000) -- (2.000,0.000);
\draw[sgrasta] (2.000,2.000) -- (3.000,1.000);
\draw[sgrasta] (2.000,2.000) -- (2.000,0.000);
\draw[sgrasta] (3.000,1.000) -- (4.000,0.000);
\draw[sgrasta] (3.000,1.000) -- (2.000,0.000);
\draw[sgrasta] (2.000,0.000) -- (4.000,0.000);
\draw[polforza] (1.000,1.590) -- (1.000,1.090) node[pos=0.05,anchor=south,inner sep=2pt] {$\fb_1$};
\draw[polforza] (2.000,2.590) -- (2.000,2.090) node[pos=0.05,anchor=south,inner sep=2pt] {$\fb_2$};
\draw[polforza] (3.000,1.590) -- (3.000,1.090) node[pos=0.05,anchor=south,inner sep=2pt] {$\fb_3$};
\draw[polestremo] (0.000,-0.840) -- (0.000,-0.090) node[pos=0.05,anchor=north,inner sep=2pt] {$\rb_A$};
\draw[polestremo] (4.000,-0.840) -- (4.000,-0.090) node[pos=0.05,anchor=north,inner sep=2pt] {$\rb_C$};
\pic at (0.000,0.000) {sgrnodo};
\pic at (2.000,0.000) {sgrnodo};
\pic at (4.000,0.000) {sgrnodo};
\pic at (1.000,1.000) {sgrnodo};
\pic at (3.000,1.000) {sgrnodo};
\pic at (2.000,2.000) {sgrnodo};
\node[anchor=east,inner sep=2pt] at (-0.080,0.000) {$A$};
\node[anchor=north,inner sep=2pt] at (2.000,-0.080) {$B$};
\node[anchor=west,inner sep=2pt] at (4.080,0.000) {$C$};
\node[anchor=south east,inner sep=2pt] at (0.970,1.030) {$D$};
\node[anchor=south west,inner sep=2pt] at (3.030,1.030) {$E$};
\node[anchor=west,inner sep=2pt] at (2.100,2.020) {$F$};
\node[polcampo,minimum size=4.3mm] at (0.220,0.980) {$0$};
\node[polcampo,minimum size=4.3mm] at (1.450,1.880) {$1$};
\node[polcampo,minimum size=4.3mm] at (2.550,1.880) {$2$};
\node[polcampo,minimum size=4.3mm] at (3.780,0.980) {$3$};
\node[polcampo,minimum size=4.3mm] at (2.000,-0.620) {$m$};
\node[polcampo,minimum size=4.3mm] at (1.000,0.360) {$a$};
\node[polcampo,minimum size=4.3mm] at (1.620,1.000) {$b$};
\node[polcampo,minimum size=4.3mm] at (2.380,1.000) {$c$};
\node[polcampo,minimum size=4.3mm] at (3.000,0.360) {$d$};
\draw[sgrrot] ([shift={(140:0.24)}]3.85,2.15) arc (140:-130:0.24);
\node at (2.0,-1.45) {(a)};
\end{scope}
\begin{scope}[shift={(9.9,3.25)},x=1.7cm,y=1.7cm]
\draw[sgrpuntone] (0.000,0.000) -- (-1.500,-1.500);
\draw[sgrtirante] (-1.500,-1.500) -- (0.000,-1.500);
\draw[sgrpuntone] (0.000,-1.000) -- (-1.000,-2.000);
\draw[sgrpuntone] (-1.000,-2.000) -- (-1.500,-1.500);
\draw[sgrpuntone] (0.000,-2.000) -- (-1.000,-1.000);
\draw[sgrtirante] (-1.000,-1.000) -- (-1.000,-2.000);
\draw[sgrpuntone] (0.000,-3.000) -- (-1.500,-1.500);
\draw[sgrpuntone] (-1.500,-1.500) -- (-1.000,-1.000);
\draw[polforza] (0.000,0.000) -- (0.000,-1.000);
\draw[polforza] (0.000,-1.000) -- (0.000,-2.000);
\draw[polforza] (0.000,-2.000) -- (0.000,-3.000);
\draw[polcostruzione] (0.000,0.000) -- (0.720,0.000) (0.000,-1.500) -- (0.720,-1.500) (0.000,-3.000) -- (0.720,-3.000);
\draw[polestremo] (0.620,-3.000) -- (0.620,-1.500) node[midway,anchor=west,inner sep=3pt] {$\rb_C$};
\draw[polestremo] (0.620,-1.500) -- (0.620,0.000) node[midway,anchor=west,inner sep=3pt] {$\rb_A$};
\pic at (0.000,0.000) {polpunto};
\pic at (0.000,-1.000) {polpunto};
\pic at (0.000,-2.000) {polpunto};
\pic at (0.000,-3.000) {polpunto};
\pic at (-1.500,-1.500) {polpunto};
\pic at (-1.000,-2.000) {polpunto};
\pic at (-1.000,-1.000) {polpunto};
\draw[fill=white,line width=.6pt] (0.000,-1.500) circle (1.5pt);
\node[anchor=south,inner sep=3.5pt] at (0.000,0.000) {$0$};
\node[anchor=north,inner sep=3.5pt] at (0.000,-3.000) {$3$};
\node[anchor=south west,inner sep=1.5pt] at (0.000,-1.000) {$1$};
\node[anchor=north west,inner sep=1.5pt] at (0.000,-2.000) {$2$};
\node[anchor=west,fill=white,inner sep=1pt,xshift=4pt] at (0.000,-1.500) {$m$};
\node[anchor=east,inner sep=4pt] at (-1.500,-1.500) {$a \equiv d$};
\node[anchor=north east,inner sep=2pt] at (-1.000,-2.000) {$b$};
\node[anchor=south east,inner sep=2pt] at (-1.000,-1.000) {$c$};
\node at (-0.6,-3.021) {(b)};
\end{scope}
\end{tikzpicture}%
